% =================================================================
%  bilevel.tex
%  Bilevel scheduling model for the 2026 FIFA World Cup group stage
%  CORS OR Challenge 2026 -- Team UofT Champions
%
%  Upper level: FIFA chooses the schedule (match -> slot, stadium).
%  Lower level: each team chooses a base camp from a discrete
%               approved set to minimise its own round-trip travel.
%  Base-camp-dependent KPIs are evaluated at the team's best response.
% =================================================================
\documentclass[12pt, letterpaper]{article}

\usepackage[top=1in, bottom=1in, left=1in, right=1in]{geometry}
\usepackage{amsmath}
\usepackage{amssymb}
\usepackage{mathtools}
\usepackage[T1]{fontenc}
\usepackage[expansion=false]{microtype}
\usepackage{setspace}
\usepackage{booktabs}
\usepackage{array}
\usepackage{enumitem}
\usepackage[hidelinks, breaklinks=true]{hyperref}

\onehalfspacing
\setlength{\parskip}{4pt}

\newcommand{\R}{\mathbb{R}}
\newcommand{\bld}[1]{\boldsymbol{#1}}
\DeclareMathOperator*{\minimise}{minimise}
\DeclareMathOperator*{\maximise}{maximise}

\title{\textbf{A Bilevel Scheduling Framework for the 2026 FIFA World
Cup Group Stage}\\[4pt]
\large with Rational Base-Camp Selection}
\author{Team UofT Champions \\ \small University of Toronto}
\date{May 2026}

\begin{document}
\maketitle

\section{Overview and Modelling Logic}

FIFA fixes the schedule before any team selects a base camp; teams then
choose a base camp \emph{in response} to the schedule they are given.
We model this exactly as a \textbf{bilevel program}:

\begin{itemize}[leftmargin=1.4em, itemsep=2pt]
  \item \textbf{Upper level (leader, FIFA).} Choose the assignment of
        the 72 matches to (time slot, stadium) pairs, subject to the
        hard feasibility constraints of the released format. The
        objective is a weighted sum of \emph{soft} KPI penalties.
  \item \textbf{Lower level (followers, the 48 teams).} Given its three
        venues and dates --- which are \emph{determined by the upper-level
        schedule} --- each team independently selects one base camp from a
        predetermined approved set $B_i$ so as to minimise its own
        round-trip travel.
\end{itemize}

The base-camp-dependent KPIs (1.2, 1.3, 1.4, 2.4) are evaluated at the
team's optimal response $b_i^\star(S)$, never at a base camp the team
would not actually choose. Because the lower-level feasible set is
\emph{finite and small} ($|B_i|$ candidate cities), we do not need KKT
conditions or a continuous-bilevel reformulation. We replace each
lower-level problem with (i) a selection constraint, and (ii) an
\emph{optimality} (no-better-alternative) constraint that forces the
chosen camp to be a minimiser. This yields a single-level mixed-integer
linear program (MILP), exact for the discrete lower level.

\bigskip
The remaining base-camp-free KPIs in the chosen set (1.6, 1.7, 2.2, 3.3,
4.1, 4.2, 5.1, 5.2, 5.3) enter the upper-level objective directly.

% =================================================================
\section{Sets, Indices, Parameters}

\subsection*{Sets and indices}
\begin{tabular}{@{}lp{11.2cm}@{}}
\toprule
$M$ & group-stage matches, $m\in M$, $|M|=72$ \\
$T$ & time slots (date$\times$time), $t\in T$; $\tau_t$ is the
      chronological time of $t$ in hours \\
$S$ & host stadiums, $s\in S$, $|S|=16$ \\
$I$ & national teams, $i\in I$, $|I|=48$ \\
$G$ & groups, $g\in G$, $|G|=12$ \\
$M_i\subset M$ & matches involving team $i$ ($|M_i|=3$) \\
$M_g\subset M$ & matches of group $g$ ($|M_g|=6$) \\
$M_g^{f}\subset M_g$ & the two final-round matches of group $g$ \\
$T_r\subset T$ & slots of round $r\in\{1,2,3\}$ \\
$S_c\subset S$ & stadiums in host country $c\in\{\mathrm{USA,MEX,CAN}\}$ \\
$B_i$ & approved base-camp cities for team $i$ (discrete); index $b$ \\
$\Omega=\bigcup_i B_i$ & set of distinct base-camp facilities; index $b$ \\
$\Omega^{\mathrm{US}}\subseteq\Omega$ & US base-camp facilities \\
$F^{\mathrm{ban}}\subseteq I$ & teams under a full US entry ban \\
$F^{\mathrm{bond}}\subseteq I$ & teams under the US visa-bond programme \\
$R\subseteq S$ & roofed / climate-controlled stadiums \\
$S^{\mathrm{US}}\subseteq S$ & US stadiums \\
$\mathcal{R}$ & key broadcast markets, index $\rho$ \\
\bottomrule
\end{tabular}

\subsection*{Parameters}
\begin{tabular}{@{}lp{11.2cm}@{}}
\toprule
$\mathrm{dist}(b,s)$ & great-circle distance, base camp $b$ to stadium
                       $s$ (km) \\
$\mathrm{tz}(s),\,\mathrm{tz}(b)$ & UTC offset of stadium / base camp (h) \\
$\mathrm{elev}(s),\,\mathrm{elev}(b)$ & elevation of stadium / base camp (m) \\
$\mathrm{cl}(s)\in\{\mathrm{W,C,E}\}$ & FIFA cluster of stadium $s$ \\
$\mathrm{ctry}(s)$ & host country of stadium $s$ \\
$\theta_t$ & local kickoff clock time of slot $t$, hours in $[0,24)$ \\
$W_{mt}$ & expected WBGT of match $m$ if played in slot $t$
           (precomputed from venue climatology; venue is implied by $x$) \\
$\sigma_i\in\{0,0.5,1\}$ & US entry/visa severity weight of team $i$
                           (KPI 1.7) \\
$E_i$ & FIFA/Elo rating of team $i$ (lower $=$ stronger) \\
$\mu_m=(E_{i(m)}^{-1}+E_{j(m)}^{-1})/2$ & match value of $m$ \\
$q_t=\sum_{\rho}\mathrm{pop}_\rho\,w_\rho(t)$ & broadcast slot quality of
                       slot $t$ \\
$\mathrm{pop}_\rho,\,w_\rho(t)$ & audience weight and prime-time overlap
                       of market $\rho$ in slot $t$ \\
$P\subseteq M$ & high-profile matches \\
$\pi^{\mathrm{bond}}$ & US visa-bond base-camp penalty (km-equiv., KPI 1.7) \\
$R_{\min}=72$, duration $105/60$ h & rest / match-length constants \\
$N_c$ & mandated match count in country $c$ (52 / 10 / 10) \\
\bottomrule
\end{tabular}

% =================================================================
\section{Decision Variables}

\paragraph{Upper level (schedule).}
\begin{align}
x_{mts} &=
  \begin{cases} 1 & m \text{ in slot } t,\text{ stadium } s\\ 0 &
  \text{else}\end{cases}
  &&\forall m\in M,\,t\in T,\,s\in S,\\
y_{gt} &=
  \begin{cases} 1 & t \text{ is the final slot of group } g\\ 0 &
  \text{else}\end{cases}
  &&\forall g\in G,\,t\in T.
\end{align}

\paragraph{Lower level (base-camp response).}
\begin{align}
z_{ib} &=
  \begin{cases} 1 & \text{team } i \text{ camps at } b\in B_i\\ 0 &
  \text{else}\end{cases}
  &&\forall i\in I,\,b\in B_i.
\end{align}

\paragraph{Auxiliary / linearization variables.}
introduced where needed below; their meaning is stated at the point of
use.

% =================================================================
\section{Hard Constraints (from \texttt{MIP.tex})}

These are imposed exactly as in the team's feasibility model; we restate
them in the present notation.

\paragraph{H1 -- Each match scheduled once.}
\begin{equation}
\sum_{t\in T}\sum_{s\in S} x_{mts}=1 \qquad\forall m\in M.
\end{equation}

\paragraph{H2 -- Round-robin incidence.}
\begin{equation}
\sum_{m\in M_i}\sum_{t\in T}\sum_{s\in S} x_{mts}=3 \qquad\forall i\in I.
\end{equation}

\paragraph{H3 -- One match per team per round.}
\begin{equation}
\sum_{m\in M_i}\sum_{t\in T_r}\sum_{s\in S} x_{mts}=1
  \qquad\forall i\in I,\;r\in\{1,2,3\}.
\end{equation}

\paragraph{H4 -- Stadium turnover.} A stadium hosts at most one match in
any rolling window of length $\Delta_{\min}+105/60$:
\begin{equation}
\sum_{m\in M}\;\sum_{\substack{t'\in T\\ \tau_t\le\tau_{t'}<\tau_t+\Delta_{\min}+105/60}}
  x_{mt's}\le 1 \qquad\forall s\in S,\,t\in T.
\end{equation}

\paragraph{H5 -- Minimum team rest (72 h kickoff-to-kickoff $+$ duration).}
\begin{equation}
\sum_{s} x_{mts}+\sum_{s} x_{m't's}\le 1
  \quad
  \begin{aligned}[t]
  &\forall i\in I,\;\forall m,m'\in M_i,\,m\neq m',\\
  &\forall t,t'\in T:\,|\tau_t-\tau_{t'}|<R_{\min}+105/60.
  \end{aligned}
\end{equation}

\paragraph{H6 -- Host-nation venue restriction.}
\begin{equation}
\sum_{t\in T}\sum_{s\in S\setminus S_i} x_{mts}=0
  \qquad\forall i\in H,\,m\in M_i.
\end{equation}

\paragraph{H7 -- Simultaneous final matches.}
\begin{align}
\sum_{t\in T} y_{gt}&=1 &&\forall g\in G,\\
\sum_{m\in M_g}\sum_{s} x_{mts}&\ge 2\,y_{gt} &&\forall g\in G,\,t\in T,\\
\sum_{m\in M_g}\sum_{s}\sum_{\substack{t'\in T\\ \tau_{t'}>\tau_t}} x_{mt's}
  &\le 2\,(1-y_{gt}) &&\forall g\in G,\,t\in T.
\end{align}

\paragraph{H8 -- Match allocation by host country.}
\begin{equation}
\sum_{m\in M}\sum_{t\in T}\sum_{s\in S_c} x_{mts}=N_c
  \qquad\forall c\in\{\mathrm{USA,MEX,CAN}\}.
\end{equation}

% =================================================================
\section{The Lower Level and its Exact Linearization}
\label{sec:lower}

\subsection{Lower-level problem}
Given the schedule, team $i$'s three venues are the stadiums on which its
matches are played. The team chooses $b\in B_i$ to minimise its perceived
round-trip travel cost (KPI 1.1), augmented by a US visa-bond penalty
defined below:
\begin{equation}
b_i^\star \in \arg\min_{b\in B_i}\;
  C_{ib}, \qquad
  C_{ib}\;=\;\underbrace{2\sum_{m\in M_i}\sum_{t\in T}\sum_{s\in S}
            \mathrm{dist}(b,s)\,x_{mts}}_{\text{round-trip travel}}
            \;+\; P_{ib}.
\label{eq:llcost}
\end{equation}
The travel term is \emph{linear} in $x$ (distances are constants; only
the venue indicator $x$ is a variable); $P_{ib}$ is a constant
visa-restriction penalty (\S\ref{sec:visa}), so the difficulty is not the
cost expression but the $\arg\min$.

\subsection{US entry restrictions on base-camp eligibility}
\label{sec:visa}
A team subject to US entry restrictions would not realistically site its
base camp in the United States. We distinguish the two restriction tiers
of KPI 1.7. Let $\Omega^{\mathrm{US}}\subseteq\Omega$ be the US
facilities.

\paragraph{Full-ban teams ($F^{\mathrm{ban}}$): hard exclusion.}
A full entry suspension forecloses a US base entirely. For every
$i\in F^{\mathrm{ban}}$ we forbid US camps:
\begin{equation}
z_{ib}=0 \qquad\forall i\in F^{\mathrm{ban}},\;
  b\in B_i\cap\Omega^{\mathrm{US}}.
\label{eq:banexcl}
\end{equation}
(Equivalently, US facilities are removed from $B_i$ in preprocessing; the
constraint form is stated for transparency.) These teams are then
\emph{required} to camp at a non-US facility, so $B_i$ must retain at
least one Canadian or Mexican option for feasibility.

\paragraph{Visa-bond teams ($F^{\mathrm{bond}}$): soft penalty.}
Players and staff are exempt from the bond, so a US base remains
physically possible but carries real financial and logistical friction.
We discourage --- rather than forbid --- US camps for these teams via a
constant additive penalty $\pi^{\mathrm{bond}}>0$ (in km-equivalent
units, so it is commensurate with the travel term):
\begin{equation}
P_{ib}=
  \begin{cases}
  \pi^{\mathrm{bond}} & \text{if } i\in F^{\mathrm{bond}}
    \text{ and } b\in\Omega^{\mathrm{US}},\\
  0 & \text{otherwise.}
  \end{cases}
\label{eq:bondpen}
\end{equation}
Because $P_{ib}$ is a \emph{constant}, the penalised cost $C_{ib}$ in
\eqref{eq:llcost} stays linear in $x$ and requires no new linearization;
it enters the realised cost \eqref{eq:realcost} and the optimality cut
\eqref{eq:condopt} unchanged, so the team genuinely weighs the bond
friction when choosing. Setting $\pi^{\mathrm{bond}}$ larger than the
maximum intra-set travel spread recovers the hard-exclusion behaviour as
a limiting case; a moderate value lets a bond team still choose a US camp
when the travel saving is large enough to outweigh the friction.

\subsection{Selection and optimality constraints}
We require exactly one camp per team:
\begin{align}
\sum_{b\in B_i} z_{ib} &= 1 &&\forall i\in I,\label{eq:selone}\\
C_i^\star &= \sum_{b\in B_i} z_{ib}\, C_{ib}
  &&\forall i\in I. \label{eq:realcost}
\end{align}

\subsection{Camp exclusivity (one team per facility)}
\label{sec:exclusive}
FIFA-approved base camps are physically exclusive: a facility hosts at
most one team. Let $\Omega=\bigcup_i B_i$ be the set of distinct
facilities. For each facility we impose
\begin{equation}
\sum_{i\,:\,b\in B_i} z_{ib} \;\le\; 1 \qquad\forall b\in\Omega.
\label{eq:exclusive}
\end{equation}

\paragraph{Consequence: the optimality cut must be made conditional.}
With exclusivity, the lower level is \emph{no longer} 48 independent
problems --- it is a joint assignment of teams to a shared, congestible
resource. A team can therefore be denied its unconstrained
travel-minimiser when another team has claimed it. A \emph{hard}
optimality cut $C_i^\star\le C_{ib}\ \forall b$ would then be
\textbf{infeasible} precisely whenever two teams' minimisers coincide.
We instead require optimality only among camps \emph{not taken by another
team} (a no-justified-envy / stability condition):
\begin{equation}
C_i^\star \;\le\; C_{ib} \;+\;
  \mathcal{M}\!\!\sum_{j\neq i\,:\,b\in B_j} z_{jb}
  \qquad\forall i\in I,\;b\in B_i,
\label{eq:condopt}
\end{equation}
where $\mathcal{M}$ is a valid big-$M$. If camp $b$ is unclaimed by any
other team, the sum is $0$ and \eqref{eq:condopt} reduces to the ordinary
cut $C_i^\star\le C_{ib}$, forcing team $i$ to beat $b$. If some team
$j\neq i$ took $b$, the sum is $1$, the right-hand side becomes
$C_{ib}+\mathcal{M}$, and the cut is deactivated --- team $i$ is not
required to beat a camp it could not have. Together
\eqref{eq:selone}--\eqref{eq:realcost}, \eqref{eq:exclusive} and
\eqref{eq:condopt} make $z$ a \emph{stable} assignment: every team holds
the cheapest camp still available to it given the others' choices.

\paragraph{Choosing $\mathcal{M}$.} Any upper bound on a realised travel
cost works; the tightest valid constant is
\begin{equation}
\mathcal{M} \;=\; \max_{i\in I}\Bigl(\max_{b\in B_i}C^{\mathrm{ub}}_{ib}
  - \min_{b\in B_i}C^{\mathrm{lb}}_{ib}\Bigr),
\end{equation}
where $C^{\mathrm{ub}}_{ib}=2\sum_{k=1}^{3}\max_{s}\mathrm{dist}(b,s)$ and
$C^{\mathrm{lb}}_{ib}=2\sum_{k=1}^{3}\min_{s}\mathrm{dist}(b,s)$ bound the
realised cost over the team's three matches; both are constants, so
$\mathcal{M}$ is precomputed. Using the difference rather than a raw cost
keeps $\mathcal{M}$ as small as possible, tightening the LP relaxation.

\paragraph{Modelling consequences.} Two points to record explicitly.
(i)~\emph{Separability is lost}: because \eqref{eq:exclusive} couples the
teams, the lower level can no longer be precomputed as a per-team table
lookup --- the camp assignment is now solved jointly with the schedule
inside the MILP. (ii)~Ties are still broken arbitrarily and remain
harmless, since all KPIs depend on the chosen camp's coordinates, not on
$z$ itself.

\subsection{Linearizing the products in $C_{ib}$ and $C_i^\star$}
\label{sec:bilinear}

Two bilinear obstacles remain.

\paragraph{(a) The travel part of $C_{ib}$ is linear in $x$ already.}
Define the \emph{travel contribution} parameter
$D_{ibs}=2\,\mathrm{dist}(b,s)$. Then
\begin{equation}
C_{ib}=\sum_{m\in M_i}\sum_{t\in T}\sum_{s\in S} D_{ibs}\,x_{mts}
       \;+\; P_{ib},
\qquad D_{ibs}=2\,\mathrm{dist}(b,s),
\end{equation}
which is linear ($P_{ib}$ is a constant); no work needed. Constraint
\eqref{eq:condopt} is therefore already linear in $x$.

\paragraph{(b) $C_i^\star=\sum_b z_{ib}C_{ib}$ is a product
$z_{ib}\cdot x_{mts}$ (binary $\times$ binary).} We replace each product
by an auxiliary variable
\begin{equation}
u_{ibmts}=z_{ib}\,x_{mts}\in\{0,1\},
\end{equation}
linearized by the standard McCormick envelope for a product of two
binaries:
\begin{equation}
u_{ibmts}\le z_{ib},\quad
u_{ibmts}\le x_{mts},\quad
u_{ibmts}\ge z_{ib}+x_{mts}-1,\quad
u_{ibmts}\ge 0.
\label{eq:mccormick}
\end{equation}
Then
\begin{equation}
C_i^\star=\sum_{b\in B_i}\sum_{m\in M_i}\sum_{t\in T}\sum_{s\in S}
   D_{ibs}\,u_{ibmts}
   \;+\;\sum_{b\in B_i} P_{ib}\,z_{ib}.
\end{equation}
The second sum is linear in $z$ ($P_{ib}$ constant), so the visa-bond
penalty needs no product linearization.
Because $z_{ib}$ and $x_{mts}$ are binary, \eqref{eq:mccormick} is
\emph{exact} (the McCormick envelope is tight at binary endpoints), so no
relaxation gap is introduced. To control model size we only generate
$u_{ibmts}$ for $m\in M_i$ (team $i$'s own three matches), so the count is
$\sum_i 3\,|B_i|\,|T|\,|S|$, which is modest.

\medskip
\emph{Remark (cheaper alternative).} The whole lower level can instead be
\textbf{precomputed}: for each team and each candidate venue-triple,
$\arg\min_b C_{ib}$ is a table lookup. Embedding it as
\eqref{eq:selone}--\eqref{eq:condopt} is the in-model version we use so that
base-camp KPIs feed back into the optimisation rather than being scored
only after the fact.

% =================================================================
\section{Base-Camp-Dependent Soft KPIs}
\label{sec:bc-kpi}

All four are evaluated at $b_i^\star$ via $z$. To express a quantity ``at
the chosen camp'' linearly, we use the same product-linearization device:
for any camp-indexed constant $\kappa_{ib}$,
\begin{equation}
\widehat{\kappa}_i \;=\; \sum_{b\in B_i} z_{ib}\,\kappa_{ib}
\label{eq:atchosen}
\end{equation}
is linear in $z$ because $\kappa_{ib}$ is constant. Quantities that also
depend on the venue (hence on $x$) reuse $u_{ibmts}$.

\subsection{KPI 1.2 -- Intra-Group Travel Dispersion}
We use the realised total travel $C_i^\star$ (already defined). For group
$g$ define the within-group range via $\max$/$\min$ epigraph variables:
\begin{align}
C^{\max}_g \ge C_i^\star,\quad
C^{\min}_g \le C_i^\star &&\forall g\in G,\,i\in g,\\
\Delta_g &= C^{\max}_g - C^{\min}_g &&\forall g\in G.
\end{align}
$\Delta_g\ge0$ automatically. (Linearizing $\max$/$\min$ by epigraph
bounds is exact because we minimise $\sum_g\Delta_g$, which pushes
$C^{\max}_g$ down and $C^{\min}_g$ up to the true extremes.) We use the
\emph{range} form rather than the coefficient of variation because CV is a
nonlinear ratio (see \S\ref{sec:nonlinkpi}). Penalty term:
$\mathrm{KPI}_{1.2}=\sum_{g} \Delta_g$.

\subsection{KPI 1.3 -- Circadian Shift Cost (Jet-Lag)}
The perceived kickoff time depends on the venue's time-zone offset
relative to the base camp, and the penalty $\phi(\cdot)$ is piecewise
linear. We precompute, for every (team $i$, camp $b$, slot $t$, stadium
$s$) the constant penalty
\begin{equation}
\Phi_{ibts}=\phi\!\Bigl(\bigl(\theta_t-[\mathrm{tz}(s)-\mathrm{tz}(b)]\bigr)\bmod 24\Bigr),
\end{equation}
since $\theta_t,\mathrm{tz}(s),\mathrm{tz}(b)$ are all parameters --- so the
nonlinear mod-24 and the piecewise $\phi$ are evaluated \emph{offline} and
enter as a constant. The team's jet-lag load is then
\begin{equation}
\mathrm{JL}_i=\sum_{b\in B_i}\sum_{m\in M_i}\sum_{t\in T}\sum_{s\in S}
  \Phi_{ibts}\,u_{ibmts},
\end{equation}
linear via the same $u$ variables. Penalty:
$\mathrm{KPI}_{1.3}=\sum_i \mathrm{JL}_i$ (or a minimax variant, below).

\subsection{KPI 1.4 -- Match-Venue Geographic Dispersion}
Cluster diversity $\mathrm{CC}_i$ does not require a base camp, but the
border-crossing count does. We use the border-crossing form because it is
the base-camp-dependent variant. Precompute the constant
\begin{equation}
\beta_{ibs}=\mathbf 1\!\left[\mathrm{ctry}(s)\neq\mathrm{ctry}(b)\right],
\end{equation}
giving
\begin{equation}
\mathrm{BC}_i=\sum_{b\in B_i}\sum_{m\in M_i}\sum_{t\in T}\sum_{s\in S}
  \beta_{ibs}\,u_{ibmts}\in\{0,1,2,3\},
\end{equation}
again linear through $u$. Penalty: $\mathrm{KPI}_{1.4}=\sum_i \mathrm{BC}_i$.

\subsection{KPI 2.4 -- Altitude Disruption Index}
The penalty $f(\Delta e)=\max(0,\Delta e-500)/1000$ depends only on the
\emph{constant} elevations of $b$ and $s$, so precompute
\begin{equation}
A_{ibs}=f\!\bigl(|\mathrm{elev}(s)-\mathrm{elev}(b)|\bigr),
\end{equation}
and set
\begin{equation}
\mathrm{ALT}_i=\sum_{b\in B_i}\sum_{m\in M_i}\sum_{t\in T}\sum_{s\in S}
  A_{ibs}\,u_{ibmts}.
\end{equation}
The $\max(0,\cdot)$ inside $f$ is resolved offline (it is a function of
constants), so no in-model $\max$ linearization is needed. Penalty
(minimax form): introduce $\mathrm{ALT}^{\max}\ge \mathrm{ALT}_i\ \forall
i$, penalise $\mathrm{ALT}^{\max}$.

% =================================================================
\section{Base-Camp-Free Soft KPIs}
\label{sec:free-kpi}

\subsection{KPI 1.6 -- Rest Asymmetry Between Opponents}
Let $r_{mi}$ be the rest (in slots/days) of team $i$ before match $m$;
this is a linear function of the slot variables. For match $m$ between
$i,j$ we need $|r_{mi}-r_{mj}|$. Linearize the absolute value with a
nonnegative epigraph variable:
\begin{align}
\mathrm{RD}_m &\ge r_{mi}-r_{mj}, &
\mathrm{RD}_m &\ge r_{mj}-r_{mi}, &
\mathrm{RD}_m&\ge 0.
\end{align}
Minimising $\sum_m\mathrm{RD}_m$ pulls $\mathrm{RD}_m$ down to the true
absolute difference, so the relaxation is exact at optimum. For the
\emph{count} objective $|\{m:\mathrm{RD}_m>0\}|$ introduce a binary
$\delta_m$ with $\mathrm{RD}_m\le \bar R\,\delta_m$ ($\bar R$ a valid
upper bound on the rest gap) and penalise $\sum_m\delta_m$. Penalty:
$\mathrm{KPI}_{1.6}=\sum_m \mathrm{RD}_m$.

\subsection{KPI 1.7 -- Entry \& Visa Restriction Exposure (ERI)}
For match $m$ at stadium $s$, exposure occurs iff $s\in S^{\mathrm{US}}$,
weighted by $\max(\sigma_{i(m)},\sigma_{j(m)})$ which is a \emph{constant}
$\bar\sigma_m=\max(\sigma_{i(m)},\sigma_{j(m)})$ known from the draw.
Hence
\begin{equation}
\mathrm{ERI}=\sum_{m\in M}\bar\sigma_m
  \sum_{t\in T}\sum_{s\in S^{\mathrm{US}}} x_{mts},
\end{equation}
fully linear (the $\max$ is over constants, resolved offline). Penalty:
$\mathrm{KPI}_{1.7}=\mathrm{ERI}$.

\subsection{KPI 2.2 -- Per-Team Heat Load}
$\mathrm{WBGT}$ at match $m$ depends on venue, date, and slot, all encoded
by $x$; $W_{mt}$ (the WBGT if $m$ is played in slot $t$, at the stadium
implied by $x$) is precomputed per (match, slot, stadium). The excess
over $28^\circ$C uses a nonnegative epigraph variable $h_m$:
\begin{align}
h_m &\ge \sum_{t,s} W_{mts}\,x_{mts} - 28, & h_m &\ge 0,
\end{align}
so $h_m=\max(0,\mathrm{WBGT}_m-28)$ at optimum (minimising pushes $h_m$
down to the bound). Team heat load $\mathrm{HL}_i=\sum_{m\in M_i} h_m$,
and we penalise the minimax $\mathrm{HL}^{\max}\ge\mathrm{HL}_i$.

\subsection{KPI 3.3 -- Round-Order Balance (First-Mover)}
For each group $g$ and non-simultaneous round $r\in\{1,2\}$, the two
matches are ordered by slot. Let $p_{ir}\in\{0,1\}$ indicate team $i$
plays the \emph{earlier} match of round $r$. Within round $r$ of group
$g$ exactly one of the two matches is earlier; define $p_{ir}$ via the
slot variables: for the two matches $m_1,m_2\in M_g$ of round $r$,
$p_{ir}=\sum_{t,s}\mathbf 1[\text{$m$ is the earlier slot}]\,x_{mts}$ for
the match $m\ni i$. The positional score is $a_i=\sum_{r\in\{1,2\}}
p_{ir}\in\{0,1,2\}$. We minimise the spread of $a_i$ via epigraph
variables $a^{\max}\ge a_i\ge a^{\min}$ and penalise $a^{\max}-a^{\min}$
(a linear surrogate for the standard deviation $\mathrm{FMB}$; see
\S\ref{sec:nonlinkpi}). ``Earlier match'' indicators are themselves
linearized with order binaries $o_{m m'}\in\{0,1\}$ and big-$M$ slot
comparisons:
\begin{equation}
\tau_{t}-\tau_{t'} \le \mathcal M\,o_{mm'},\qquad
\tau_{t'}-\tau_{t} \le \mathcal M\,(1-o_{mm'}),
\end{equation}
for the slots $t,t'$ chosen for $m,m'$ (themselves products with $x$,
linearized as in \S\ref{sec:bilinear} if needed; in practice the
round/slot structure fixes orderings combinatorially and $o$ can be read
off without products).

\subsection{KPI 4.1 -- Venue-Load Balance}
Let $n_s=\sum_{m,t} x_{mts}$ be the match count at stadium $s$. We
minimise the load \emph{range} (linear surrogate for the CV):
\begin{align}
n^{\max}\ge n_s,\quad n^{\min}\le n_s &&\forall s\in S,\qquad
\mathrm{KPI}_{4.1}=n^{\max}-n^{\min}.
\end{align}

\subsection{KPI 4.2 -- Same-City Overlap \& Fan Accessibility}
Same-city-same-day overlap is a product $x_{mts}x_{m't's}$-type term over
matches sharing a city and date. Introduce $w_{mm'}\in\{0,1\}$ with
\begin{equation}
w_{mm'}\ge \xi_{m}+\xi_{m'}-1,\quad w_{mm'}\le \xi_m,\quad w_{mm'}\le
\xi_{m'},
\end{equation}
where $\xi_m=\sum_{t\in T_d}\sum_{s\in S_{\text{city}}} x_{mts}$ indicates
$m$ is in the given city on the given day (a linear sum of $x$). Then
$\mathrm{SCO}=\sum_{\text{city},\,d}\sum_{m<m'} w_{mm'}$. Penalty:
$\mathrm{KPI}_{4.2}=\mathrm{SCO}$.

\subsection{KPI 5.1 -- Prime-Time Alignment}
Slot quality $q_t$ is a constant. Maximising alignment:
\begin{equation}
\mathrm{PT}=\sum_{m\in M}\sum_{t\in T}\sum_{s\in S} q_t\, x_{mts},
\end{equation}
linear. Enter as $-\mathrm{PT}$ in the (minimisation) objective.

\subsection{KPI 5.2 -- Marquee-Match Slot Quality \& Overlap}
$\mathrm{MSQ}=\sum_{m}\mu_m\sum_{t,s} q_t x_{mts}$ (linear, $\mu_m$
constant). Simultaneous high-profile penalty over $P$: for
$m,m'\in P$, $\mathrm{OVL}\!:\ \pi_{mm'}\in\{0,1\}$ with
$\pi_{mm'}\ge \zeta_{mt}+\zeta_{m't}-1$ summed over shared slots, where
$\zeta_{mt}=\sum_s x_{mts}$. Net term: $-(\mathrm{MSQ}-\lambda\,
\mathrm{OVL})$.

\subsection{KPI 5.3 -- Host-City Economic Equity}
Venue commercial value $\mathrm{VC}_s=\sum_{m,t}\mu_m q_t x_{mts}$ is
linear. The Gini coefficient is nonlinear (\S\ref{sec:nonlinkpi}); we use
the \emph{mean absolute difference} surrogate, which shares Gini's
ordering and is linearizable:
\begin{equation}
\mathrm{MAD}=\sum_{s<s'} \eta_{ss'},\quad
\eta_{ss'}\ge \mathrm{VC}_s-\mathrm{VC}_{s'},\;
\eta_{ss'}\ge \mathrm{VC}_{s'}-\mathrm{VC}_s,\;\eta_{ss'}\ge0,
\end{equation}
and maximise equity by penalising $\mathrm{MAD}$. (Gini $=
\mathrm{MAD}/(2\bar{\mathrm{VC}}\,|S|)$; with $\sum_s\mathrm{VC}_s$ fixed
by the schedule's total value, minimising MAD aligns with maximising
$\mathrm{HCE}$.)

% =================================================================
\section{Objective Function}

All penalties are normalised to a common scale via precomputed bounds
$\overline{K}$ (the maximum attainable value of each KPI, used as a
divisor). The leader solves:
\begin{align}
\minimise_{x,y,z,u,\dots}\quad
& \lambda_{1.2}\frac{\sum_g\Delta_g}{\overline{K}_{1.2}}
+ \lambda_{1.3}\frac{\sum_i\mathrm{JL}_i}{\overline{K}_{1.3}}
+ \lambda_{1.4}\frac{\sum_i\mathrm{BC}_i}{\overline{K}_{1.4}}
+ \lambda_{1.6}\frac{\sum_m\mathrm{RD}_m}{\overline{K}_{1.6}}
\notag\\
& + \lambda_{1.7}\frac{\mathrm{ERI}}{\overline{K}_{1.7}}
+ \lambda_{2.2}\frac{\mathrm{HL}^{\max}}{\overline{K}_{2.2}}
+ \lambda_{2.4}\frac{\mathrm{ALT}^{\max}}{\overline{K}_{2.4}}
+ \lambda_{3.3}\frac{a^{\max}-a^{\min}}{\overline{K}_{3.3}}
\notag\\
& + \lambda_{4.1}\frac{n^{\max}-n^{\min}}{\overline{K}_{4.1}}
+ \lambda_{4.2}\frac{\mathrm{SCO}}{\overline{K}_{4.2}}
- \lambda_{5.1}\frac{\mathrm{PT}}{\overline{K}_{5.1}}
- \lambda_{5.2}\frac{\mathrm{MSQ}-\lambda\,\mathrm{OVL}}{\overline{K}_{5.2}}
+ \lambda_{5.3}\frac{\mathrm{MAD}}{\overline{K}_{5.3}}
\\[4pt]
\text{s.t.}\quad
& \text{H1--H8 (hard feasibility),}\notag\\
& \text{\eqref{eq:selone}--\eqref{eq:condopt}, \eqref{eq:exclusive} (lower-level best response),}
  \notag\\
& \text{\eqref{eq:mccormick} (product linearization) and all KPI
  epigraph / indicator constraints above,}\notag\\
& x,y,z,u,\delta,o,w,\pi\in\{0,1\},\quad \text{epigraph vars}\ \ge 0.
\notag
\end{align}
Maximised KPIs (5.1, 5.2) enter with a negative sign. The weights
$\lambda_\bullet\ge0$ are set by the analyst; a $\sum\lambda=1$
normalisation lets the same model trace a Pareto frontier by sweeping the
weights.

% =================================================================
\section{Summary of Linearizations}
\label{sec:nonlinkpi}

\begin{center}
\renewcommand{\arraystretch}{1.25}
\begin{tabular}{@{}p{2.0cm} p{4.6cm} p{6.4cm}@{}}
\toprule
\textbf{KPI} & \textbf{Nonlinearity} & \textbf{Linearization} \\
\midrule
Lower level & discrete $\arg\min_b$ with shared-camp exclusivity &
  selection \eqref{eq:selone} $+$ exclusivity \eqref{eq:exclusive} $+$
  conditional (big-$M$) optimality cut \eqref{eq:condopt}; yields a
  stable assignment \\
US entry restr.\ (1.7) & ban $=$ hard exclusion; bond $=$ friction &
  fix $z_{ib}=0$ for $F^{\mathrm{ban}}$ on US camps; constant penalty
  $\pi^{\mathrm{bond}}$ added to $C_{ib}$ for $F^{\mathrm{bond}}$ (stays
  linear) \\
$C_i^\star$, 1.3, 1.4, 2.4 & binary product $z_{ib}x_{mts}$ &
  McCormick \eqref{eq:mccormick}; exact at binary endpoints \\
1.2 & range $\max-\min$; CV ratio &
  epigraph vars for $\max,\min$; \emph{range} used in place of CV \\
1.3 / 2.4 & mod-24, piecewise $\phi$, $\max(0,\cdot)$ in $f$ &
  evaluated \emph{offline} into constants $\Phi_{ibts}$, $A_{ibs}$ \\
1.6 & absolute value $|r_{mi}-r_{mj}|$ &
  two-sided epigraph $\mathrm{RD}_m\ge\pm(\cdot)$; binary $\delta_m$ for
  the count form \\
1.7 / 5.1 / 5.2 (MSQ) & $\max$ over constants; products with constants &
  resolved offline; remain linear in $x$ \\
2.2 & $\max(0,\mathrm{WBGT}-28)$ & nonneg.\ epigraph $h_m$ \\
3.3 & ordering ``earlier match''; std.\ dev.\ $\mathrm{FMB}$ &
  big-$M$ order binaries $o_{mm'}$; range $a^{\max}-a^{\min}$ surrogate
  for std.\ dev. \\
4.1 & CV of venue loads & range $n^{\max}-n^{\min}$ surrogate \\
4.2 / 5.2 (OVL) & binary products (same city/day; same slot) &
  McCormick on $w_{mm'},\pi_{mm'}$ \\
5.3 & Gini coefficient &
  mean-absolute-difference surrogate $\mathrm{MAD}$ with two-sided
  epigraph $\eta_{ss'}$ \\
\bottomrule
\end{tabular}
\end{center}

\paragraph{Why surrogates for CV / std.\ dev.\ / Gini.}
The coefficient of variation (1.2, 4.1), the standard deviation (3.3),
and the Gini coefficient (5.3) are all \emph{ratios or square roots} and
are not MILP-representable without nonlinear (e.g.\ SOCP) machinery.
Their \emph{range} and \emph{mean-absolute-difference} counterparts are
piecewise-linear, share the same monotone ``more equal $\Rightarrow$
smaller'' direction, and are standard fairness surrogates in the
scheduling literature. If exact CV/Gini are required, the model becomes a
MISOCP; we recommend the linear surrogates for tractability and report
the exact CV/Gini \emph{post hoc} on the optimised schedule.

\end{document}